\section{Introduction}
\label{sec:intro}

The \emph{Binary Symmetric Channel} (BSC) is the canonical model for digital
transmission over a noisy medium: each transmitted bit is independently flipped
with probability $p$ \cite{shannon1948}. Shannon's channel coding theorem
establishes the capacity
\begin{equation}
  C(p) = 1 - H_b(p) = 1 + p\log_2 p + (1-p)\log_2(1-p),
\end{equation}
as a hard upper bound on reliable communication rate.  Classical \emph{Error
Correcting Codes} (ECCs) approach this bound through algebraic structure:
Hamming codes correct single-bit errors within a block \cite{hamming1950},
Reed-Solomon codes correct byte-level bursts \cite{reed_solomon1960}, and
turbo/LDPC codes push to within a fraction of a dB of capacity
\cite{berrou1993,gallager1962}.

A fundamentally different correction mechanism has emerged with large language
models (LLMs): rather than exploiting algebraic redundancy, LLMs exploit the
statistical structure of natural language to \emph{infer} the most probable
original message from a corrupted version.  This linguistic redundancy operates
at a semantic level that is orthogonal to the bit-level algebra of ECCs.

\textbf{Contributions.}  This paper presents the first systematic benchmark
comparing classical ECC methods against LLM-based correction on the BSC:
\begin{enumerate}
  \item A unified experiment framework that sweeps $p \in (0, 0.5)$ and
        measures bit error rate (BER), character error rate (CER), word error
        rate (WER), and BLEU score for each method.
  \item An empirical analysis showing which methods approach the Shannon bound
        and at what noise levels they fail.
  \item Evidence that LLMs recover semantic content at noise levels where ECCs
        degrade, without violating Shannon's theorem (LLMs exploit a \emph{source
        model}, not additional channel capacity).
  \item A combined ECC+LLM pipeline that outperforms either method alone across
        all tested values of $p$.
\end{enumerate}
